\section{Конструкторский раздел}

В данном разделе приводится описание структуры разрабатываемой базы данных. В
соответствии с аналитическим разделом была выбрана гибридная архитектура хранения данных.
Для реляционной базы данных каждой выделенной сущности ставится в соответствие таблица, а
для графовой базы данных описываются структуры узлов и ребер. Рассматриваются
спроектированные ограничения целостности, а также ролевая модель управления доступом на
уровне приложения.

\subsection{Проектирование базы данных}

На основе формализации данных, проведенной в аналитическом разделе, в разрабатываемой
реляционной базе данных выделены следующие таблицы:
\begin{itemize}
    \item teams --- таблица команд разработчиков;
    \item users --- таблица пользователей приложения;
    \item invites --- таблица приглашений для регистрации;
    \item it\_systems --- таблица систем.
\end{itemize}

В таблицах \ref{tbl:teams} --- \ref{tbl:it_system} определены типы данных и ограничения
для каждого столбца перечисленных реляционных таблиц.

\begin{center}
    \begin{tabularx}{\textwidth}{|X|X|X|X|}
        \caption{\label{tbl:teams}Информация о таблице teams} \\ \hline
        Столбец & Значение & Тип & Ограничение \\ \hline
        \endfirsthead

        \multicolumn{4}{l}{{Продолжение таблицы\ \thetable{}}} \\ \hline
        Столбец & Значение & Тип & Ограничение \\ \hline
        \endhead

        \endfoot
        \endlastfoot

        id & Идентификатор & UUID & Не null, первичный ключ \\ \hline
        name & Название & Строка & Не null, уникальный \\ \hline
        description & Описание & Строка & \\ \hline
    \end{tabularx}
\end{center}

\begin{center}
    \begin{tabularx}{\textwidth}{|X|X|X|X|}
        \caption{\label{tbl:user_account}Информация о таблице users} \\ \hline
        Столбец & Значение & Тип & Ограничение \\ \hline
        \endfirsthead

        \multicolumn{4}{l}{{Продолжение таблицы\ \thetable{}}} \\ \hline
        Столбец & Значение & Тип & Ограничение \\ \hline
        \endhead

        \endfoot
        \endlastfoot

        id & Идентификатор & UUID & Не null, первичный ключ \\ \hline
        username & Имя пользователя & Строка & Не null, уникальный \\ \hline
        email & Электронная почта & Строка & Не null, уникальный \\ \hline
        password\_hash & Хэш пароля & Строка & Не null \\ \hline
        role & Идентификатор роли & Целое число & Не null \\ \hline
        team\_id & Идентификатор команды & UUID & Внешний ключ \\ \hline
    \end{tabularx}
\end{center}

\begin{center}
    \begin{tabularx}{\textwidth}{|X|X|X|X|}
        \caption{\label{tbl:invites}Информация о таблице invites} \\ \hline
        Столбец & Значение & Тип & Ограничение \\ \hline
        \endfirsthead

        \multicolumn{4}{l}{{Продолжение таблицы\ \thetable{}}} \\ \hline
        Столбец & Значение & Тип & Ограничение \\ \hline
        \endhead

        \endfoot
        \endlastfoot

        id & Идентификатор & UUID & Не null, первичный ключ \\ \hline
        status & Статус приглашения & Целое число & Не null \\ \hline
        target\_email & Электронная почта получателя приглашения & Строка & Не null, уникальный \\ \hline
        role & Роль нового пользователя & Целое число & Не null \\ \hline
        code & Уникальный код & Строка & Не null, уникальный \\ \hline
        expiration\_date & Дата истечения & Дата и время & Не null \\ \hline
        team\_id & Идентификатор команды & UUID & Не null, внешний ключ \\ \hline
        activated\_by\_uid & Идентификатор нового пользователя & UUID & Не null, внешний ключ \\ \hline
    \end{tabularx}
\end{center}

\begin{center}
    \begin{tabularx}{\textwidth}{|X|X|X|X|}
        \caption{\label{tbl:it_system}Информация о таблице it\_systems} \\ \hline
        Столбец & Значение & Тип & Ограничение \\ \hline
        \endfirsthead

        \multicolumn{4}{l}{{Продолжение таблицы\ \thetable{}}} \\ \hline
        Столбец & Значение & Тип & Ограничение \\ \hline
        \endhead

        \endfoot
        \endlastfoot

        id & Идентификатор & UUID & Не null, первичный ключ \\ \hline
        name & Название & Строка & Не null \\ \hline
        description & Описание & Строка & \\ \hline
        created\_at & Дата создания & Дата и время & Не null \\ \hline
        updated\_at & Дата обновления & Дата и время & Не null \\ \hline
        team\_id & Идентификатор команды & UUID & Не null, внешний ключ \\ \hline
    \end{tabularx}
\end{center}

Для хранения топологии инфраструктуры в графовой базе данных спроектирована структура
узлов и связей. Данные сущности не имеют жесткой табличной схемы, но подчиняются правилам
модели графа свойств.

В графовой базе данных определен узел Component, атрибуты которого представлены в таблице
\ref{tbl:graph_node}.

\begin{center}
    \begin{tabularx}{\textwidth}{|X|X|X|X|}
        \caption{\label{tbl:graph_node}Информация об атрибутах узла Component} \\ \hline
        Атрибут & Значение & Тип & Ограничение \\ \hline
        \endfirsthead

        \multicolumn{4}{l}{{Продолжение таблицы\ \thetable{}}} \\ \hline
        Атрибут & Значение & Тип & Ограничение \\ \hline
        \endhead

        \endfoot
        \endlastfoot

        id & Идентификатор & UUID & Не null, уникальный \\ \hline
        name & Название & Строка & Не null \\ \hline
        type & Тип компонента & Число & Не null \\ \hline
        description & Описание & Строка & \\ \hline
        system\_id & Идентификатор системы & UUID & Не null \\ \hline
    \end{tabularx}
\end{center}

Связь между узлами определяется направленным ребром DEPENDS\_ON. Информация об атрибутах
данного ребра представлена в таблице \ref{tbl:graph_edge}.

\begin{center}
    \begin{tabularx}{\textwidth}{|X|X|X|X|}
        \caption{\label{tbl:graph_edge}Информация об атрибутах связи DEPENDS\_ON} \\ \hline
        Атрибут & Значение & Тип & Ограничение \\ \hline
        \endfirsthead

        \multicolumn{4}{l}{{Продолжение таблицы\ \thetable{}}} \\ \hline
        Атрибут & Значение & Тип & Ограничение \\ \hline
        \endhead

        \endfoot
        \endlastfoot

        id & Идентификатор & UUID & Не null, уникальный \\ \hline
        protocol & Тип протокола & Число & Не null \\ \hline
        severity & Уровень критичности & Число & Не null \\ \hline
    \end{tabularx}
\end{center}

На рисунке \ref{img:relational-dbml.pdf} представлена полная схема разрабатываемой реляционной базы данных.

\jimg{relational-dbml.pdf}{Схема реляционной БД в нотации DBML}

\clearpage

\subsection{Ограничения целостности данных}

\subsubsection{Целостность таблиц и узлов}
Для обеспечения целостности реляционных таблиц каждая строка имеет уникальный
идентификатор, являющийся первичным ключом. Первичные ключи назначены полям id во всех
спроектированных таблицах. Для графовой базы данных на уровне СУБД устанавливается
ограничение уникальности на свойство id для всех узлов с меткой Component, что
предотвращает дублирование узлов графа.

\subsubsection{Целостность полей}
Для корректного функционирования системы накладываются ограничения на значения отдельных
полей. В таблице users поля username и email должны быть уникальными для
предотвращения создания дублирующихся учетных записей. В таблице invites поля code и 
target\_email также являются уникальным, а значение поля expiration\_date должно быть 
строго больше даты создания приглашения. В таблице it\_systems дата создания не может превышать текущее
системное время. В графовой базе данных атрибут type узла Component может принимать
только определенные значения: микросервис, база данных, брокер сообщений, внешний API.

\subsubsection{Целостность ссылок}
Обеспечение ссылочной целостности в реляционной базе данных реализуется с помощью внешних
ключей. Для каждого значения внешнего ключа в дочерней таблице должна существовать запись
с соответствующим первичным ключом в родительской таблице. В таблице \ref{tbl:fk}
представлены все внешние ключи реляционной схемы.

\begin{center}
    \begin{tabularx}{\textwidth}{|X|X|X|}
        \caption{\label{tbl:fk}Внешние ключи таблиц реляционной базы данных} \\ \hline
        Таблица & Внешний ключ & Ссылка на родительскую таблицу \\ \hline
        \endfirsthead

        \multicolumn{3}{l}{{Продолжение таблицы\ \thetable{}}} \\ \hline
        Таблица & Внешний ключ & Ссылка на родительскую таблицу \\ \hline
        \endhead

        \endfoot
        \endlastfoot

        users & teams\_id & teams (id) \\ \hline
        invites & teams\_id & teams (id) \\ \hline
        invites & activated\_by\_uid & users (id) \\ \hline
        it\_systems & teams\_id & teams (id) \\ \hline
    \end{tabularx}
\end{center}

\subsection{Ролевая модель на уровне приложения}
Для обеспечения безопасности и разграничения доступа к информации реализована ролевая
модель управления доступом на программном уровне. Базы данных не содержат индивидуальных
учетных записей для каждого пользователя системы. Вместо этого приложение подключается к
хранилищам данных с использованием единых учетных записей, а вся логика
проверки прав доступа и фильтрации запросов выполняется на уровне контроллеров и сервисов
разрабатываемого приложения.

В системе определены уровни доступа для пользователей, описанные в таблице~\ref{tbl:roles-desc}

\begin{center}
    \begin{tabularx}{\textwidth}{|X|X|}
        \caption{\label{tbl:roles-desc}Описание ролей системы} \\ \hline
        Роль & Права \\ \hline
        \endfirsthead

        \multicolumn{2}{l}{{Продолжение таблицы\ \thetable{}}} \\ \hline
        Роль & Права \\ \hline
        \endhead

        \hline
        \endfoot
        \endlastfoot

        Незарегистрированный пользователь & Имеет доступ исключительно к функциям
        регистрации, авторизации и валидации кода приглашения. \\ \hline
        
        Разработчик & Обладает правами на чтение данных о пользователях своей
        команды, а также на просмотр архитектурных схем систем, привязанных к его
        команде. Пользователь с данной ролью может инициировать запросы на визуализацию
        топологии, но не имеет прав на внесение изменений в структуру графа. \\ \hline

        SRE & Наследует все права разработчика, а также получает доступ к
        аналитическим функциям приложения. Данная роль позволяет запускать
        алгоритмические расчеты каскадных отказов, моделировать сценарии сбоев
        компонентов и запрашивать вычисление критических путей и единых точек отказа в графе. \\ \hline

        Архитектор & Обладает расширенными правами, включающими в себя возможности
        SRE-инженера, а также права на изменение топологии. Архитектор может отправлять
        запросы на создание, обновление и удаление систем, добавление новых узлов и
        настройку связей между ними в графовой базе данных. \\ \hline

        Администратор & Обладает полными правами в системе управления доступом.
        Пользователь с этой ролью может просматривать список всех зарегистрированных
        пользователей, создавать и удалять команды, генерировать коды приглашений. \\ \hline
    \end{tabularx}
\end{center}

\subsection[Проверка нахождения разработанной реляционной БД в 3НФ]{Проверка нахождения разработанной \\реляционной БД в 3НФ}
Для подтверждения корректности спроектированной схемы реляционной базы данных проведена 
проверка на соответствие требованиям третьей нормальной формы. Проверка выполняется 
последовательно путем анализа функциональных зависимостей атрибутов в каждой таблице.

Первая нормальная форма требует атомарности всех атрибутов и наличия первичного ключа 
для каждого отношения. В спроектированной базе данных все поля содержат элементарные значения: 
идентификаторы UUID, строковые данные, целые числа и временные метки. Использование массивов 
или вложенных структур в рамках одной ячейки таблицы не допускается. Каждое отношение 
обладает уникальным первичным ключом id. На основе этих признаков сделан вывод о соответствии 
базы данных требованиям первой нормальной формы.

Вторая нормальная форма подразумевает нахождение отношения в первой нормальной форме и отсутствие 
частичных функциональных зависимостей неключевых атрибутов от первичного ключа. Частичная 
зависимость возможна только в случае использования составного первичного ключа. В разработанной 
схеме первичные ключи всех таблиц являются простыми и состоят из одного столбца id. Следовательно, 
любой неключевой атрибут таблицы может зависеть только от первичного ключа целиком. Таким 
образом, база данных соответствует критериям второй нормальной формы.

Третья нормальная форма требует выполнения условий второй нормальной формы и отсутствия 
транзитивных зависимостей неключевых атрибутов от первичного ключа. Это означает, что 
ни один неключевой атрибут не должен определять значения другого неключевого атрибута 
внутри одного отношения.

В таблице users неключевые атрибуты username, email, password\_hash, role и team\_id зависят 
только от первичного ключа id. Значение роли пользователя или идентификатор его 
команды не позволяют однозначно определить его почтовый адрес или хэш пароля.

В таблице invites неключевые атрибуты status, target\_email, role, code, expiration\_date, 
team\_id и activated\_by\_uid зависят толькое от первичного ключа id. Транзитивные зависимости 
между неключевыми полями отсутствуют, так как атрибуты описывают свойства конкретного 
приглашения и не связаны между собой функциональными отношениями.

В таблице it\_systems неключевые атрибуты name, description, created\_at, updated\_at 
и team\_id определяются идентификатором системы. Описание системы или дата ее создания 
не зависят от идентификатора команды или названия системы.

В таблице teams атрибуты name и description характеризуют непосредственно команду, 
определяемую первичным ключом.

На основе проведенного анализа функциональных зависимостей установлено, 
что в спроектированной схеме базы данных каждый неключевой атрибут является фактом о 
первичном ключе и ни о чем другом. Таким образом, разработанная реляционная база 
данных находится в третьей нормальной форме.

\numberlesssubsection{Выводы по конструкторскому разделу}
В данном разделе была спроектирована структура хранения данных для системы анализа
отказоустойчивости. Подробно описаны таблицы реляционной базы данных и структуры графовой
базы данных, определены связи между ними и ограничения целостности. Формализована ролевая
модель управления доступом на уровне приложения, обеспечивающая безопасность
взаимодействия с системой и разграничение функциональных возможностей пользователей.
